%% secoes/01b-trabalhos-relacionados.tex — incluído por main.tex

%%=====================================================================%%
\section{Trabalhos relacionados}\label{sec:relacionados}
%%=====================================================================%%

No Brasil, os nascidos vivos são registrados no Sistema de Informações sobre Nascidos Vivos (SINASC). 
A vinculação desse sistema a registros socioeconômicos, como o Cadastro Único e o Programa Bolsa Família, deu origem a coortes de base populacional, entre elas a Coorte de 100 Milhões de Brasileiros utilizada nesta pesquisa, que permitem investigar esses desfechos em famílias de baixa renda \cite{pita2018}.

Para esta seção, foram buscados no PubMed e no Google Scholar artigos publicados entre 2020 e 2026 que utilizassem o SINASC ou a Coorte de 100 Milhões para investigar a relação entre fatores sociais e econômicos e ao menos um dos quatro desfechos analisados neste trabalho: baixo peso ao nascer, nascimento pré-termo, malformação congênita e Apgar baixo no quinto minuto. 
Cada artigo foi conferido quanto à base de dados, ao desfecho e ao período antes de ser incluído. 

A Tabela~\ref{tab:rel-trabalhos} resume os $19$ artigos encontrados na busca, agrupados por desfecho e descritos pela base de dados, recorte populacional, técnica de investigação e os fatores identificados como de maior risco.

%% Tabela-resumo (aparece logo após o parágrafo que a cita)
%% secoes/tabela-trabalhos-relacionados.tex — incluído por 01b-trabalhos-relacionados.tex
%%
%% Tabela-resumo dos trabalhos relacionados (2020-2026), em página
%% retrato, com quebra automática entre páginas (longtable).
%% O ano aparece na própria citação; o desfecho é o título de cada
%% grupo; o DOI de cada artigo está na lista de referências.
%% Para incluir um artigo novo, copie uma linha do grupo certo e
%% acrescente a referência em referencias/referencias.bib.

\begin{table}[H]
\centering
\caption{Trabalhos de 2020 a 2026 que investigaram fatores sociais e econômicos associados aos desfechos deste estudo a partir do SINASC ou da Coorte de 100 Milhões de Brasileiros.}
\label{tab:rel-trabalhos}

\scriptsize
\setlength{\tabcolsep}{3.5pt}
\renewcommand{\arraystretch}{1.04}

% Distribuição otimizada: colunas curtas com largura exata e 'Y' (coluna X alinhada à esquerda)
% absorvendo todo o espaço restante para a coluna de fatores de risco
\begin{tabularx}{\textwidth}{@{}p{2.3cm} p{1.85cm} p{2.55cm} p{2.35cm} Y@{}}
\toprule
\textbf{Estudo} & \textbf{Base de dados} & \textbf{Recorte} & \textbf{Técnica} & \textbf{Fatores de maior risco} \\
\midrule
\multicolumn{5}{@{}l}{\textbf{Baixo peso ao nascer}} \\
\cite{falcao2020} & 100M + SINASC & Brasil, baixa renda; 8,8 mi nascidos a termo (2001--2015) & Regressão logística & Poucas consultas de pré-natal, primiparidade, baixa escolaridade, sexo feminino, idade $\geq$35 anos, mãe preta \\
\cite{falcao2021} & 100M + SINASC & Brasil, baixa renda; 5,5 mi (2012--2015) & Regressão logística multinomial & PIG: sem pré-natal, analfabetismo, mãe preta, parda ou indígena, sem companheiro, idade 14--20 ou 35--49 anos \\
\cite{falcao2023} & CIDACS + BF & Brasil, baixa renda; 4,3 mi (2012--2015) & Escore de propensão & BF reduziu o BPN em 11\%, mais entre indígenas, mães com $\leq$3 anos de estudo e com $<$7 consultas \\
\cite{almeida2026} & SINASC & São Paulo; 2,5 mi (2019--2023) & Aprendizado de máquina & Gestação múltipla, $\leq$6 consultas, anomalia, idade $\leq$19 ou $\geq$35 anos (descritivo) \\

\addlinespace[3.5pt]
\multicolumn{5}{@{}l}{\textbf{Nascimento pré-termo}} \\
\cite{rocha2022} & CIDACS & Brasil, multíparas de baixa renda; 3,5 mi (2001--2015) & Modelos de transição (reg. logística) & Até 3 consultas, intervalo entre partos $<$12 meses, idade $\geq$35; só no incidente: indígena, 14--19 anos, adensamento \\
\cite{alberton2023} & SINASC & Brasil (2011--2021) & Séries temporais & Indígenas, 1--3 anos de estudo, idades extremas, 1--3 consultas, região Norte \\
\cite{ortelan2024} & 100M + BF + SINASC & Brasil, multíparas de baixa renda; 1 mi (2012--2015) & Escore de propensão & BF reduziu em 31\% o pré-termo extremo; sem efeito no total \\
\cite{victor2025} & SINASC & Brasil, por UF (2014--2023) & Séries temporais & Sem escolaridade, indígenas, pretas, idades extremas, Norte e Nordeste \\

\addlinespace[3.5pt]
\multicolumn{5}{@{}l}{\textbf{Malformação congênita}} \\
\cite{trevilato2022} & SINASC & Rio Grande do Sul; 5.250 casos (2012--2015) & Caso-controle; reg. logística & Idade $>$40 anos, nenhuma consulta, $<$4 anos de estudo, mãe preta, perda fetal anterior \\
\cite{anele2022} & SINASC & Porto Alegre, por nível de desenv. (2000--2017) & Regressão de Poisson & Idade materna avançada; $<$8 anos de estudo só nas áreas mais desenvolvidas \\
\cite{paixao2022} & SINASC + reg. eventos & Brasil; 11,4 mi (2015--2018) & Regressão logística & Síndrome do Zika: mãe preta ou parda, solteira, $<$12 anos de estudo, menos consultas \\
\cite{fernandes2025} & SINASC + SIM & Brasil; 26,1 mi (2012--2020) & Reg. logística hierárquica & Idade $>$40 anos, gestação múltipla, sem pré-natal, mãe preta \\
\cite{melkan2025} & SINASC & Brasil; 29 mi (2011--2020) & Descritiva e tendência temporal & Mães pretas, gestação múltipla, sexo masculino; tubo neural inverso à escolaridade \\

\addlinespace[3.5pt]
\multicolumn{5}{@{}l}{\textbf{Apgar baixo}} \\
\cite{magalhaes2023} & SINASC & Brasil; 9 mi (1999 e 2018--2019) & Regressão logística & Prematuridade, baixo peso, anomalia; idade $\geq$35, $<$8 anos de estudo, solteira, $<$6 consultas, RN preto \\
\cite{pinto2024} & SINASC & Brasil, meninas de 10--14 anos (2011--2021) & Ecológico; análise espacial & Mais Apgar baixo e BPN entre meninas de 10--14 anos, maioria negra, Norte e Nordeste \\
\cite{lima2025} & SINASC & Indígenas do Território Wayamu (PA); 713 nasc. (2013--2022) & Descritiva & Cerca de 9\% com Apgar $<$8; gravidez na adolescência, poucas consultas \\
\cite{delcarretto2025} & SINASC + temperatura & São Paulo; 3,2 mi de baixo risco (2013--2019) & Caso-cruzado & Calor extremo, com efeito em mães com $<$12 anos de estudo e pardas \\

\addlinespace[3.5pt]
\multicolumn{5}{@{}l}{\textbf{Múltiplos desfechos} (pré-termo, BPN, anomalia e/ou Apgar)} \\
\cite{reboucas2024} & SINASC + SIM & Brasil; 21,3 mi (2012--2019) & Frações atribuíveis & Mães indígenas e pretas, sobretudo as de menor escolaridade \\
\cite{pescarini2024} & CIDACS + SIM & Brasil, mães migrantes; 10,3 mi (2011--2018) & Regressão logística & Migrantes: mais anomalia (sobretudo pretas), menos pré-termo e BPN; Apgar sem diferença \\
\bottomrule
\end{tabularx}

\vspace{3pt}
\parbox{\textwidth}{\raggedright\fontsize{6.5}{7.5}\selectfont
\textbf{Fonte:} elaborada pelos autores. 100M: Coorte de 100 Milhões de Brasileiros; CIDACS: Coorte de Nascimentos do CIDACS (SINASC vinculado à 100M); SIM: Sistema de Informações sobre Mortalidade; BPN: baixo peso ao nascer; PIG: pequeno para a idade gestacional; BF: Bolsa Família; RN: recém-nascido; UF: unidade federativa; mi: milhões.}
\end{table}

%%---------------------------------------------------------------------%%
\subsection{Fatores associados aos desfechos neonatais adversos}\label{subsec:rel-desfechos}
%%---------------------------------------------------------------------%%

Para o baixo peso ao nascer, os estudos sobre a Coorte de 100 Milhões, que analisaram o baixo peso entre nascidos a termo e o nascimento pequeno para a idade gestacional, apontam como fatores de maior risco o número reduzido de consultas de pré-natal, a baixa escolaridade materna, a raça/cor preta, a primiparidade e as idades maternas extremas \cite{falcao2020,falcao2021}. Dois trabalhos tratam o desfecho como efeito de uma política ou de uma desigualdade: o recebimento do Bolsa Família durante a gestação reduziu em 11\% a chance de baixo peso, sobretudo entre mães indígenas e de menor escolaridade \cite{falcao2023}, e 7,2\% dos casos seriam evitados se todas as mulheres tivessem as condições das mulheres brancas \cite{reboucas2024}.

Para o nascimento pré-termo, as séries temporais nacionais do SINASC registram as maiores prevalências entre mães indígenas e pretas, com baixa escolaridade, nas idades extremas, com poucas consultas de pré-natal e residentes no Norte e no Nordeste \cite{alberton2023,victor2025}. Na Coorte de 100 Milhões, a adolescência e o adensamento domiciliar elevaram o risco apenas da prematuridade incidente, e não da recorrente \cite{rocha2022}, e o Bolsa Família reduziu em 31\% os pré-termo extremos, sem efeito sobre o total \cite{ortelan2024}. Entre mães indígenas, 22,2\% dos nascimentos pré-termo seriam evitáveis na ausência de desigualdade étnico-racial \cite{reboucas2024}.

Para a malformação congênita, os estudos, entre eles um dedicado à síndrome congênita do Zika, associam o desfecho à raça/cor preta, à ausência ou ao número reduzido de consultas de pré-natal, à idade materna avançada e à baixa escolaridade \cite{fernandes2025,trevilato2022,paixao2022,melkan2025}. O peso desses fatores, porém, varia conforme o recorte: em Porto Alegre, a baixa escolaridade só se associou ao desfecho nas áreas mais desenvolvidas \cite{anele2022}, e filhos de mães migrantes internacionais tiveram maior chance de anomalia, concentrada sobretudo entre mães pretas \cite{pescarini2024}.

O Apgar baixo é o desfecho menos estudado sob a perspectiva social. Apenas um trabalho o analisa por regressão multivariada, associando-o à idade materna a partir de 35 anos, à baixa escolaridade, à mãe solo, a menos de seis consultas de pré-natal e à raça/cor preta do recém-nascido \cite{magalhaes2023}. Nos demais, ele aparece como desfecho secundário ou em análises descritivas, como entre mães de 10 a 14 anos \cite{pinto2024} e no Território Indígena Wayamu, no Pará \cite{lima2025}, ou como efeito do calor extremo, mais evidente entre mães com menos de 12 anos de estudo e pardas \cite{delcarretto2025}.

%%---------------------------------------------------------------------%%
\subsection{Descoberta de subgrupos e aprendizado de máquina}\label{subsec:rel-subgrupos}
%%---------------------------------------------------------------------%%

Poucos trabalhos aplicam aprendizado de máquina a esses desfechos, e em geral com fins preditivos. No SINASC de São Paulo, regressão logística e floresta aleatória foram comparadas na predição do baixo peso, com os fatores associados apresentados apenas de forma descritiva \cite{almeida2026}. Um modelo de aprendizado profundo treinado sobre 22.104 pares mãe--recém-nascido prevê 24 desfechos neonatais adversos \cite{defrancesco2023}. Em ambos, o objetivo é classificar um novo caso, e não descrever de forma legível os grupos em que o desfecho se concentra.

A busca por subgrupos tem justamente esse objetivo. Em Guiné-Bissau, a combinação de análise espacial, regressão e aprendizado de máquina identificou grupos de crianças com mortalidade elevada, que se mantiveram em períodos distintos e podem orientar intervenções dirigidas \cite{rieckmann2024}. Na Coorte de 100 Milhões, a técnica já foi aplicada ao baixo peso ao nascer \cite{cunha2024}, mas não aos demais desfechos neonatais adversos.

%%---------------------------------------------------------------------%%
\subsection{Síntese e posicionamento deste trabalho}\label{subsec:rel-lacuna}
%%---------------------------------------------------------------------%%

Os trabalhos revisados apresentam três limitações em relação ao objetivo desta pesquisa. A primeira é a técnica: regressões, séries temporais e estatísticas descritivas estimam o efeito de cada atributo isoladamente ou descrevem prevalências, mas não retornam as combinações de características em que o desfecho se concentra. A segunda é o tratamento dos grupos vulneráveis, que em geral entram como variáveis de ajuste ou como estratos de uma mesma análise, e não como recortes analisados separadamente, cada um com a sua própria prevalência de referência. A terceira é a cobertura dos desfechos: a malformação congênita e, sobretudo, o Apgar baixo foram pouco investigados, e apenas dois trabalhos analisam vários desfechos na mesma população \cite{reboucas2024,pescarini2024}, ainda assim modelando cada um separadamente.

Este trabalho parte desse precedente \cite{cunha2024} e o estende de um para quatro desfechos, acrescentando o nascimento pré-termo, a malformação congênita e o Apgar baixo, o que permite comparar quais características se repetem entre eles e quais são próprias de cada um.
